\documentclass[12pt,a4paper]{ctexart}
\usepackage{amsmath}
\usepackage{amsfonts}
\usepackage{amssymb}
\usepackage[margin=1in]{geometry}

\usepackage{titlesec}

\usepackage[hidelinks]{hyperref}

\numberwithin{equation}{section}
\titleformat{\section}[block]{\bfseries\Large}{\thesection\quad}{0pt}{}

\begin{document}
\begin{center}
\thispagestyle{empty}

\vspace{2.5cm}
{\fontsize{44}{48}\bfseries\selectfont Notes on Szabo}

\vfill
\begin{center}
\Large  \textit{by} \textbf{Shangjia Liu} 
\end{center} 
{\Large 2026}
\vspace{1.75cm}
\end{center}

\newpage
\pagenumbering{roman}  
\setcounter{page}{2}  
\tableofcontents

\newpage
\pagenumbering{arabic}
\setcounter{page}{1}
\section{Many-Electron Wave Functions and Operators}
{本章介绍量子化学的基本概念、技巧及必要的记号。具体有三方面:
多电子算符 (如哈密顿算符)的结构和多电子波函数的形式(即 Slater 
行列式及其线性组合);如 何求算符在 Slater 行列式之间的矩阵元;
Hartree-Fock 近似的基本思想。之后的章节就是 Hartree-Fock 
近似和以 HF 为起点的高级方法,所以现在学的知识在以后的章节中
极有用处。}


\subsection{The Electronic Problem}
本书重点关注如何近似求解非相对论不含时 Schrödinger 方程：
\begin{equation}
  \hat{H} |\Phi\rangle = E |\Phi\rangle 
  \label{eq:schrodinger}
\end{equation}
若用原子单位制, 则 N 个电子 M 个核的哈密顿量就是:

\begin{equation}
  \hat{H}=-\sum_{i=1}^N\frac{1}{2}\nabla_i^2-
  \sum_{A=1}^{M}\frac{1}{2M_A}\nabla_A^2-
  \sum_{A=1}^{M}\sum_{i=1}^N\frac{Z_A}{r_{iA}}+
  \sum_{i=1}^{N}\sum_{j>i}^N\frac{1}{r_{ij}}+
  \sum_{A=1}^M\sum_{B>A}^M\frac{Z_AZ_B}{R_{AB}} 
  \label{eq:hamiltonian}
\end{equation}
式 \eqref{eq:hamiltonian}中的第一项是全 部电子的动能算符,
 第二项是核的动能算符, 第三项是核与电子间的 Coulomb 吸引势能,
第四项和第五项分别是电子间和核之间的排斥势.

\subsubsection{Atomic Units}
本书采用原子单位制, 其中  \(\hbar = m_e = e = 1\)
(即普朗克常数、电子质量和电子电荷都取值为 1),
因此能量单位是 Hartree~($E_{\mathrm{h}}$), 长度单位是 Bohr~($a_0$)
.原子单位制下的Schrödinger方程为:
\begin{equation}
  \left( -\frac{1}{2} \nabla^{\prime 2} - \frac{1}{r'} \right) \phi' = E \phi'
  \label{eq:schrodinger-atomic}
\end{equation}
对于基态氢原子,该方程的解可以导出能量 $E = -0.5$~Hartree.
\subsubsection{The Born--Oppenheimer Approximation}
由于原子核远比电子重,所以它们运动得较慢.因此,很大程度上可将
分子中的电子视作在位置固定的核产生的势场中运动.在这个近似下,式
\eqref{eq:hamiltonian}中第二项, 即核动能项,可以被忽略;式中最后一项,
即核间排斥势可以视作常量.任意常数加到算符上产生的效果就是
本征值加一常数,而本征函数不变.式\eqref{eq:hamiltonian}中的剩余项称作电子
的哈密顿量,或描述在 M 个点电荷产生的势场中运动的 N 个电子的哈密顿量,
\begin{equation}
\hat{H}_{elec}=-\sum_{i=1}^N\frac{1}{2}\nabla_i^2-
\sum_{A=1}^{M}\sum_{i=1}^N\frac{Z_A}{r_{iA}}+
\sum_{i=1}^{N}\sum_{j>i}^N\frac{1}{r_{ij}}  
\end{equation}
电子哈密顿量的 Schrödinger 方程为
\begin{equation}
  \hat{H}_{elec} |\Phi\rangle = E_{elec} |\Phi\rangle
\end{equation}
其解就是电子波函数
\begin{equation}
  |\Phi_{\text{elec}}\rangle=\Phi_{\text{elec}}(\{\mathbf{r}_i\}; \{\mathbf{R}_A\})
\end{equation}
它描述电子的运动,并显式依赖于电子坐标,而以参数的方式依赖于核坐标,
电子能量也参数地依赖于核坐标:
\begin{equation}
  E_{\text{elec}} = E_{\text{elec}}(\{\mathbf{R}_A\})
  \label{eq:Eelec}
\end{equation}
它以参数依赖核坐标是指:对于每种不同的核构型, $\Phi_{\text{elec}}$
（作为电子坐标的函数）的形式都不一样。核坐标并不显式地出现在 $\Phi_{\text{elec}}$
中。核固定构型下的总能量还需包含核排斥能。
\begin{equation}
  E_{\text{total}} = E_{\text{elec}} +
   \sum_{A=1}^M\sum_{B>A}^M\frac{Z_AZ_B}{R_{AB}}
\end{equation}
\subsubsection{The Antisymmetry or Pauli Principle}
方程式\eqref{eq:Eelec}中的电子能量 $E_{\text{elec}}$ 是核坐标的函数,
 但它并不依赖于电子自旋.因此,我们引入两个自旋函数$\alpha(\omega)$和
 $\beta(\omega)$,它们分别对应自旋向上和自旋向下的电子.函数变量 $\omega$
 无明显意义，只需要自旋函数正交完备
\begin{equation}
  \langle \alpha | \alpha \rangle = \langle \beta | \beta \rangle = 1, \quad
  \langle \alpha | \beta \rangle = 0
\end{equation}
由于哈密顿算符中并无自旋,仅仅在波函数中强加自旋无法产生任何后果.
但若 对波函数加上一个额外的约束,就能得到一个完善的理论:
 多电子波函数在交换任意 两个电子坐标的操作下,必须是反对称的
\begin{equation}
  \hat{P}_{ij} |\Phi\rangle = -|\Phi\rangle
  \label{pauli}
\end{equation}
这个额外的要求有时叫作反对称原理, 它是我们熟悉的
 Pauli 不相容原理的一种常见表述.即波函数不仅要满足
  Schrödinger 方程,还要满足反对称原理\eqref{pauli}.
满足反对称原理的波函数可以用 Slater 行列式来构造
\subsection{Orbitals, Slater Determinants, and Basis Functions}
\subsubsection{Spin Orbitals and Spatial Orbitals}
常认为空间分子轨道构成一组正交基
\begin{equation}
  \langle \psi_p | \psi_q \rangle = \delta_{pq}
\end{equation}
若空间轨道集合 $\{\psi_i\}$ 为正交归一的，则用它可以精确地展开任意函数
\begin{equation}
  f(\mathbf{r}) = \sum_i^{\infty} a_i \psi_i(\mathbf{r})
\end{equation}
实际运算不可能用到完备集，而取其中有限个空间轨道 $\{\psi_p\}$ .

为完整描述一个电子, 需要指定电子自旋.描述自旋和空间两部分的波函数叫作自旋轨道, 
$\chi(x)$, 其中 x 代表自旋坐标与空间坐标.一个空间轨道可以构造两个自旋轨道
\begin{equation}
  \chi(\mathbf{x})=\begin{cases}
    \psi(\mathbf{r})\alpha(\omega) \\
    \text{or}\\
    \psi(\mathbf{r})\beta(\omega)
  \end{cases}
\end{equation}
\subsubsection{Hartree Product}\label{sec:HartreeProduct}
无相互作用 $N$ 电子的集合，其哈密顿量形式如下
\begin{equation}
  \hat{H} = \sum_{i=1}^N \hat{h}(i)
\end{equation}
其中 $\hat{h}(i)$ 是描述第 i 个电子的单电子算符.算符 $\hat{h}(i)$ 
有自己的一组本征函数 $\{\chi_p(i)\}$ ，我们将其视作一组自旋轨道${\{\chi_p(j)\}}$
\begin{equation}
  \hat{h}(i) \chi_j(i) = \epsilon_j \chi_j(i)
\end{equation}
对于无相互作用的 $N$ 电子系统, 其波函数可以写成单电子波函数的乘积
\begin{equation}
  \Psi^{HP}(\mathbf{x}_1, \mathbf{x}_2, \ldots, \mathbf{x}_N) = \chi_i(\mathbf{x}_1) \chi_j(\mathbf{x}_2) \cdots \chi_k(\mathbf{x}_N)
\end{equation}
这个波函数叫作 Hartree 积,是$\hat{H}$的一个本征函数
\begin{equation}
  \hat{H} \Psi^{HP} =E\Psi^{HP}=
   \left( \sum_{i=1}^N \epsilon_i \right) \Psi^{HP}
\end{equation}

\vspace{1em}
\noindent\rule{\linewidth}{0.4pt}

Hartree 积是无相互作用哈密顿量的本征函数，本征值为占据轨道能量之和。

\textbf{证明：} Hartree 积为
\[
  \Psi^{HP}(\mathbf{x}_1, \mathbf{x}_2, \ldots, \mathbf{x}_N) = \chi_i(\mathbf{x}_1) \chi_j(\mathbf{x}_2) \cdots \chi_k(\mathbf{x}_N)
\]
其中 $i, j, \ldots, k$ 为任意 $N$ 个占据轨道指标。无相互作用哈密顿量 $\hat{H} = \sum_{p=1}^N \hat{h}(p)$ 作用其上：
\[
  \hat{H} \Psi^{HP} = \sum_{p=1}^N \hat{h}(p) \bigl[ \chi_i(\mathbf{x}_1) \chi_j(\mathbf{x}_2) \cdots \chi_k(\mathbf{x}_N) \bigr]
\]
$\hat{h}(p)$ 仅作用于第 $p$ 个坐标，其余因子视为常数。由 $\hat{h} \chi_q = \epsilon_q \chi_q$ 得
\[
  \hat{h}(1) \Psi^{HP} = \epsilon_i \Psi^{HP},\quad
  \hat{h}(2) \Psi^{HP} = \epsilon_j \Psi^{HP},\ \ldots,\ 
  \hat{h}(N) \Psi^{HP} = \epsilon_k \Psi^{HP}
\]
逐项相加：
\[
  \hat{H} \Psi^{HP}
  = (\epsilon_i + \epsilon_j + \cdots + \epsilon_k) \Psi^{HP}
\]
故 Hartree 积是 $\hat{H}$ 的本征函数，本征值为各占据轨道能量之和。

\noindent\rule{\linewidth}{0.4pt}
\vspace{1em}

Hartree 积是无相关波函数，因为在体积微元$d\mathbf{x}_1$中找到电子1 
的概率与在体积微元$d\mathbf{x}_2$中找到电子2的概率是独立的（不相关的）
但实际上电子间每时每刻都在通过 Coulomb（库仑）相互作用排斥，所以两电子间是
明显相互关联着的.
\subsubsection{Slater Determinant}
满足反对称原理\eqref{pauli}的波函数可以用 Slater 行列式来构造.
对于 N 个电子占据 N 个自旋轨道 $\chi_i, \chi_j, \ldots, \chi_k$ 的情况, Slater 行列式定义为
\begin{equation}
  \Psi^{SD}(\mathbf{x}_1, \mathbf{x}_2, \ldots, \mathbf{x}_N) = \frac{1}{\sqrt{N!}}
  \begin{vmatrix}
    \chi_i(\mathbf{x}_1) & \chi_j(\mathbf{x}_1) & \cdots & \chi_k(\mathbf{x}_1) \\
    \chi_i(\mathbf{x}_2) & \chi_j(\mathbf{x}_2) & \cdots & \chi_k(\mathbf{x}_2) \\
    \vdots & \vdots & \ddots & \vdots \\
    \chi_i(\mathbf{x}_N) & \chi_j(\mathbf{x}_N) & \cdots & \chi_k(\mathbf{x}_N)
  \end{vmatrix}
\end{equation}
注意Slater行列式行表示电子.交换任意两行（即交换两个电子坐标）
会改变行列式的符号，因此满足反对称原理.若两电子占据同一自旋轨道,
 则行列式中两列相同, 行列式值为零, 满足 Pauli 不相容原理.
 为简便，引入归一化的Slater行列式的简写符号：隐去归一化因子，
 仅写出行列式的对角项
\begin{equation}
\Psi^{SD}(\mathbf{x}_1, \mathbf{x}_2, \ldots, \mathbf{x}_N) 
= |\chi_i (\mathbf{x}_1) \chi_j (\mathbf{x}_2) \cdots \chi_k (\mathbf{x}_N) \rangle
=|\chi_i \chi_j \cdots \chi_k \rangle
\end{equation}
由于在Slater行列式中自旋反平行的电子的运动仍然缺失关联作用，
也习惯上称Slater行列式为无相关波函数。
\subsubsection{The Hartree--Fock Approximation}
描述 $N$ 电子系统基态最简单的反对称波函数是单Slater行列式
\begin{equation}
  |\Phi_0\rangle = |\chi_1 \chi_2 \cdots \chi_N \rangle   
\end{equation}
变分原理给出上形式的波函数，最优的波函数应该给出最低的能量
\begin{equation}
  E_0 = \min_{\{\chi_i\}} \langle \Phi_0 | \hat{H} | \Phi_0 \rangle
\end{equation}

Hartree-Fock 近似的精髓在于将复杂的多电子问题转换为单电子问题，
而单电子问题是以平均的方式处理电子之间的排斥.由于Hartree-Fock方程式是非线性
方程，只能迭代求解.这种方法称为自洽场方法(self-consistent field method).
\subsubsection{The Minimal Basis \texorpdfstring{$\mathrm{H}_2$}{H2} Model}       
现在来介绍一个简单的模型体系, 此模型贯穿全书, 我们会用它来展示许多量子化
学的思想和方法. 它就是我们熟悉的极小基 MO-LCAO 下$\mathrm{H}_2$的分子.

氢原准确的 $1s$ 轨道为
\begin{equation}
  \psi_{1s}(\mathbf{r}-\mathbf{R}) = 
  \left( \frac{\zeta^3}{\pi} \right)^{1/2} e^{-\zeta |r-R|}
\end{equation}
$\mathbf{\zeta}$是轨道指数,本书取1.0.但本书多数使用Gaussian轨道，因为它在积分时
更容易，Gaussian轨道形式如下
\begin{equation}
  \psi_{1s}(\mathbf{r}-\mathbf{R}) = 
  \left(\frac{2\alpha}{\pi}\right)^{3/4} e^{-\alpha |\mathbf{r}-\mathbf{R}|^2}
\end{equation}
利用两个定域原子轨道$\phi_1$,$\phi_2$, 可以采取线性组合的方式
构建两个离域分子轨道. 对称组合生成 gerade 对称的成键分子轨道 
\begin{equation}
  \psi_1 = \frac{1}{\sqrt{2(1+S_{12})}} (\phi_1 + \phi_2)
\end{equation}
相反地也可以构建反对称组合, 生成具有 ungerade 对称性的反键分子轨道
\begin{equation}
  \psi_2 = \frac{1}{\sqrt{2(1-S_{12})}} (\phi_1 - \phi_2)
\end{equation}

\vspace{1em}
\noindent\rule{\linewidth}{0.4pt}

分子轨道 $\psi_1$ 和 $\psi_2$ 组成一组正交基。

\textbf{证明：} 假设原子轨道 $\phi_1$ 和 $\phi_2$ 是正交归一的，即 $\langle \phi_1 | \phi_1 \rangle = 1$, $\langle \phi_2 | \phi_2 \rangle = 1$, $\langle \phi_1 | \phi_2 \rangle = S_{12}$。计算内积：
\[
\langle \psi_1 | \psi_2 \rangle = \frac{1}{\sqrt{2(1+S_{12})}} \frac{1}{\sqrt{2(1-S_{12})}} \langle \phi_1 + \phi_2 | \phi_1 - \phi_2 \rangle
\]
展开：
\[
\langle \phi_1 + \phi_2 | \phi_1 - \phi_2 \rangle = \langle \phi_1 | \phi_1 \rangle - \langle \phi_1 | \phi_2 \rangle + \langle \phi_2 | \phi_1 \rangle - \langle \phi_2 | \phi_2 \rangle = 1 - S_{12} + S_{12} - 1 = 0
\]
因此，
\[
\langle \psi_1 | \psi_2 \rangle = \frac{1}{2 \sqrt{(1+S_{12})(1-S_{12})}} \cdot 0 = 0
\]
故 $\psi_1$ 和 $\psi_2$ 正交。

\noindent\rule{\linewidth}{0.4pt}
\vspace{1em}

Hartree-Fock 基态可以简写为
\begin{equation}
  |\Phi_0\rangle = |\psi_1 \bar{\psi}_1 \rangle
\end{equation}
这种符号代表两个电子占据同一个空间轨道$\psi_1$, 一个自旋是$\alpha$，
一个是 $\beta$。从上下文可以清楚地分辨出 $\psi_1$ 是代表空间轨道
还是由 $\psi_1$ 和 $\alpha$ 自旋函数组成的自旋轨道.
\subsubsection{Excited Determinants}
单重激发态是指在 Hartree-Fock 基态中占据某轨道$\chi_a$
的电子被提升到了某虚轨道$\chi_r$上
\begin{equation}
  |\Phi_a^r\rangle = |\chi_1 \cdots \chi_{a-1} \chi_r \chi_{a+1} \cdots \chi_N \rangle
\end{equation}

双重激发态是$\chi_a$和$\chi_b$上的电子激发到虚轨道$\chi_r$和$\chi_s$上
\begin{equation}
  |\Phi_{ab}^{rs}\rangle = |\chi_1 \cdots \chi_{a-1} \chi_r \chi_{a+1} \cdots \chi_{b-1} \chi_s \chi_{b+1} \cdots \chi_N \rangle
\end{equation}
以此类推, 可以定义三重激发态, 四重激发态等.某种程度上, 行列式的重数越高,
在近似描述体系的真实态时, 其重要性就越低. 
虽然激发态行列式描述的不是真实的激发态, 但重要的是, 它们可以作为 $N$ 电子基函数
来展开精确的 $N$ 电子态.
\subsubsection{Form of the Exact Wave Function and 
Configuration Interaction}
$N$ 电子问题中的基态和激发态精确波函数都可以写成由一组完全集
$\{\chi_i(\mathbf{x})\}$所生产的全部可能的$N$电子Slater行列式的线性组合

由于所有可能的行列式都可以借助基态 Hartree-Fock 行列式 (的激发) 来表示, 
所 以可将体系任意态对应的精确波函数写作:
\begin{equation}
  |\Psi\rangle = c_0 |\Phi_0\rangle + \sum_{a,r} c_a^r |\Phi_a^r\rangle + \sum_{a<b, r<s} c_{ab}^{rs} |\Phi_{ab}^{rs}\rangle + \cdots
\end{equation}
所以无穷集合 $\{|\Psi_i\rangle\}=\{|\Psi_0\rangle, |\Psi_a^r\rangle, |\Psi_{ab}^{rs}\rangle, \ldots\}$ 
是 $N$ 电子 Hilbert 空间的一组完备基函数，并可用于展开$N$电子波函数。由完备集
$\{|\Psi_i\rangle\}$ 产生的哈密顿算符矩阵其本征值就是基态和激发态的准确能量.
由于每个$|\Psi_i\rangle$都对应一种自旋轨道的组态(组态即生成$|\Psi_i\rangle$
的自旋轨道集合),因此这种方法叫做组态相互作用(configuration interaction,CI)

极小基$H_2$的精确波函数$|\Psi_0\rangle$的对称性与Hartree-Fock近似下的基态
一样，都有$g$对称性.因此, $|\Psi_0\rangle$只能由具有$g$对称性的行列式线性组合而成.
\begin{equation}
  |\Psi_0\rangle = c_0 |\Phi_0\rangle +
   c_{1\bar{1}}^{2\bar{2}} |\Phi_{1\bar{1}}^{2\bar{2}}\rangle= 
   c_0 |\Phi_0\rangle + c_{12}^{34} |\Phi_{12}^{34}\rangle
\end{equation}
确定了上述展开式中的系数, 也就确定了精确波函数$|\Psi_0\rangle$和精确能量
$\langle \Psi_0 | \hat{H} | \Psi_0 \rangle$.为此，我们将full CI
矩阵对角化，也即将以$|\Phi_0\rangle$和$|\Phi_{1\bar{1}}^{2\bar{2}}\rangle$
为基底的$2\times 2$哈密顿矩阵对角化
\begin{equation}
 \mathbf{H}= \begin{pmatrix}
    \langle \Phi_0 | \hat{H} | \Phi_0 \rangle & \langle \Phi_0 | \hat{H} | \Phi_{1\bar{1}}^{2\bar{2}} \rangle \\
    \langle \Phi_{1\bar{1}}^{2\bar{2}} | \hat{H} | \Phi_0 \rangle & \langle \Phi_{1\bar{1}}^{2\bar{2}} | \hat{H} | \Phi_{1\bar{1}}^{2\bar{2}} \rangle
  \end{pmatrix}
  \label{eq:full-ci-2x2}
\end{equation}
\subsection{Operators and Matrix Elements}
本节来讲如何求得某算符在正交归一轨道构成的行列式之间的矩阵元.
\subsubsection{Minimal Basis \texorpdfstring{$\mathrm{H_2}$}{H2} Matrix Elements}
我们先讨论极小基$H_2$模型的full CI矩阵元\eqref{eq:full-ci-2x2}.其精确基态
是Hartree-Fock基态$|\Phi_0\rangle=|1\bar{1}\rangle$和
双重激发态$|\Phi_{1\bar{1}}^{2\bar{2}}\rangle=|2\bar{2}\rangle$的线性组合.我们需要求
对角元$\langle\Phi_0|\hat{H}|\Phi_0\rangle$
和$\langle\Phi_{1\bar{1}}^{2\bar{2}}|\hat{H}|\Phi_{1\bar{1}}^{2\bar{2}}\rangle$
以及非对角元$\langle\Phi_0|\hat{H}|\Phi_{1\bar{1}}^{2\bar{2}}\rangle$.

双电子体系的哈密顿量为
\begin{equation}
  \hat{H} = \hat{h}(1) + \hat{h}(2) + \frac{1}{r_{12}}
\end{equation}
为了方便，把总哈密顿量分为单电子和双电子部分
\begin{equation}
  \hat{O}_1 = \hat{h}(1) + \hat{h}(2)
\label{singleelec}
\end{equation}
\begin{equation}
  \hat{O}_2 = \frac{1}{r_{12}}
\end{equation}

考虑矩阵元$\langle\Phi_0|\hat{O}_1|\Phi_0\rangle$，由式\eqref{singleelec}可知
它是两项之和，不妨先计算
\begin{equation}
\begin{aligned}
  \langle\Phi_0|\hat{h}(1)|\Phi_0\rangle &=
   \int d\mathbf{x}_1 d\mathbf{x}_2\, \Phi_0^*(\mathbf{x}_1, \mathbf{x}_2)\, \hat{h}(1) \,\Phi_0(\mathbf{x}_1, \mathbf{x}_2) \\ 
  &=\frac12 \int d\mathbf{x}_1 d\mathbf{x}_2\,
   \bigl[\chi_1^*(\mathbf{x}_1)\chi_2^*(\mathbf{x}_2)-\chi_2^*(\mathbf{x}_1)\chi_1^*(\mathbf{x}_2)\bigr] \\ 
  &\qquad\times \hat{h}(1) \,
   \bigl[\chi_1(\mathbf{x}_1)\chi_2(\mathbf{x}_2)-\chi_2(\mathbf{x}_1)\chi_1(\mathbf{x}_2)\bigr] \\ 
  &=\frac12 \int d\mathbf{x}_1 d\mathbf{x}_2\,
    \Bigl\{ 
      \chi_1^*(\mathbf{x}_1)\chi_2^*(\mathbf{x}_2) \,\hat{h}(1) \,\chi_1(\mathbf{x}_1)\chi_2(\mathbf{x}_2) \\ 
  &\qquad\qquad + \chi_1^*(\mathbf{x}_1)\chi_2^*(\mathbf{x}_2) \,\hat{h}(1) \,\chi_2(\mathbf{x}_1)\chi_1(\mathbf{x}_2) \\ 
  &\qquad\qquad - \chi_2^*(\mathbf{x}_1)\chi_1^*(\mathbf{x}_2) \,\hat{h}(1) \,\chi_1(\mathbf{x}_1)\chi_2(\mathbf{x}_2) \\ 
  &\qquad\qquad + \chi_2^*(\mathbf{x}_1)\chi_1^*(\mathbf{x}_2) \,\hat{h}(1) \,\chi_2(\mathbf{x}_1)\chi_1(\mathbf{x}_2) 
    \Bigr\} .
\end{aligned}
\end{equation}
由自旋轨道正交归一性，上式中交叉项积分为零，化简得
\begin{equation}
  \langle\Psi_0|\hat{h}(1)|\Psi_0\rangle
  = \frac{1}{2} \int d\mathbf{x}_1\, \chi_1^*(\mathbf{x}_1)\hat{h}(\mathbf{r}_1)\chi_1(\mathbf{x}_1)
  + \frac{1}{2} \int d\mathbf{x}_1\, \chi_2^*(\mathbf{x}_1)\hat{h}(\mathbf{r}_1)\chi_2(\mathbf{x}_1)
\end{equation}
同样可以得到 $\langle\Psi_0|\hat{h}(2)|\Psi_0\rangle = \langle\Psi_0|\hat{h}(1)|\Psi_0\rangle$，那么
\begin{equation}
  \langle\Psi_0|\hat{O}_1|\Psi_0\rangle
  = \int d\mathbf{x}_1\, \chi_1^*(\mathbf{x}_1)\hat{h}(\mathbf{r}_1)\chi_1(\mathbf{x}_1)
  + \int d\mathbf{x}_1\, \chi_2^*(\mathbf{x}_1)\hat{h}(\mathbf{r}_1)\chi_2(\mathbf{x}_1)
\end{equation}
上式中的积分叫做单电子积分，我们引入如下记号表示包含自旋轨道的单电子积分
\begin{equation}
  \langle i|\hat{h}|j\rangle = \langle\chi_i|\hat{h}|\chi_j\rangle
  = \int d\mathbf{x}_1\, \chi_i^*(\mathbf{x}_1)\hat{h}(\mathbf{r}_1)\chi_j(\mathbf{x}_1)
\end{equation}
即有
\begin{equation}
  \langle\Psi_0|\hat{O}_1|\Psi_0\rangle
  = \langle 1|\hat{h}|1\rangle + \langle 2|\hat{h}|2\rangle
\end{equation}

 现在来计算 $\hat{O}_2$ 的矩阵元

\begin{equation}
\begin{aligned}
\langle\Psi_0|\hat{O}_2|\Psi_0\rangle
&= \int d\mathbf{x}_1 d\mathbf{x}_2\;
   \Bigl[ 2^{-1/2} \bigl( \chi_1(\mathbf{x}_1)\chi_2(\mathbf{x}_2)
                         - \chi_2(\mathbf{x}_1)\chi_1(\mathbf{x}_2) \bigr) \Bigr]^* \\
&\qquad \times r_{12}^{-1} \;
   \Bigl[ 2^{-1/2} \bigl( \chi_1(\mathbf{x}_1)\chi_2(\mathbf{x}_2)
                         - \chi_2(\mathbf{x}_1)\chi_1(\mathbf{x}_2) \bigr) \Bigr] \\[4pt]
&= \frac{1}{2} \int d\mathbf{x}_1 d\mathbf{x}_2\;
   \Bigl\{
      \chi_1^*(\mathbf{x}_1)\chi_2^*(\mathbf{x}_2)\, r_{12}^{-1}\,
      \chi_1(\mathbf{x}_1)\chi_2(\mathbf{x}_2) \\
&\qquad + \chi_2^*(\mathbf{x}_1)\chi_1^*(\mathbf{x}_2)\, r_{12}^{-1}\,
      \chi_2(\mathbf{x}_1)\chi_1(\mathbf{x}_2)
      - \chi_1^*(\mathbf{x}_1)\chi_2^*(\mathbf{x}_2)\, r_{12}^{-1}\,
      \chi_2(\mathbf{x}_1)\chi_1(\mathbf{x}_2) \\
&\qquad - \chi_2^*(\mathbf{x}_1)\chi_1^*(\mathbf{x}_2)\, r_{12}^{-1}\,
      \chi_1(\mathbf{x}_1)\chi_2(\mathbf{x}_2)
   \Bigr\}.
\end{aligned}
\label{eq:O2-matrix}
\end{equation}
由于$r_{12}=r_{21}$,上式第二项积分变量可以交换，所以该项与第一项相等.
类似的，第三项与第四项也相等.因此有
\begin{equation}
  \begin{aligned}
  \langle\Psi_0|\hat{O}_2|\Psi_0\rangle
  = \int d\mathbf{x}_1 d\mathbf{x}_2\;
    \Bigl\{&
      \chi_1^*(\mathbf{x}_1)\chi_2^*(\mathbf{x}_2)\, r_{12}^{-1}\,
      \chi_1(\mathbf{x}_1)\chi_2(\mathbf{x}_2) \\
      - &\chi_1^*(\mathbf{x}_1)\chi_2^*(\mathbf{x}_2)\, r_{12}^{-1}\,
      \chi_2(\mathbf{x}_1)\chi_1(\mathbf{x}_2)
    \Bigr\}.
  \end{aligned}
\end{equation}
上面的积分就是双电子积分之一例，得名来源于积分变量包括电子1与电子2的空间坐标以及两个自旋坐标.
习惯上常将所有双电子积分中的积分变量都写成电子1和2的坐标.现引入记号来表示含自旋轨道的双电子积分
\begin{equation}
  \langle ij|kl\rangle
  =\langle\chi_i \chi_j |\chi_k \chi_l \rangle
  = \int d\mathbf{x}_1 d\mathbf{x}_2\,
     \chi_i^*(\mathbf{x}_1)\chi_j^*(\mathbf{x}_2)\, r_{12}^{-1}\,
     \chi_k(\mathbf{x}_1)\chi_l(\mathbf{x}_2)
\end{equation}
那么就有
\begin{equation}
  \langle\Psi_0|\hat{O}_2|\Psi_0\rangle
  = \langle 12|12\rangle - \langle 12|21\rangle
\end{equation}
Hartree--Fock基态能量就成为
\begin{equation}
  E_0 = \langle 1|h|1\rangle + \langle 2|h|2\rangle + \langle 12|12\rangle - \langle 12|21\rangle
\end{equation}
\subsubsection{Notations for One- and Two-Electron Integrals}
由于双电子积分经常以如下的组合形式出现，在此引入一个特别记号来表示反对称
双电子积分
\begin{equation}
  \begin{aligned}
    \langle ij||kl\rangle
    &= \langle ij|kl\rangle - \langle ij|lk\rangle \\[2pt]
    &= \int d\mathbf{x}_1 d\mathbf{x}_2\,
       \chi_i^*(\mathbf{x}_1)\chi_j^*(\mathbf{x}_2)\,
       r_{12}^{-1}\,
       \bigl[ \chi_k(\mathbf{x}_1)\chi_l(\mathbf{x}_2)
            - \chi_l(\mathbf{x}_1)\chi_k(\mathbf{x}_2) \bigr]
  \end{aligned}
\end{equation}
自旋轨道
\begin{equation}
  \left[i|h|j\right]=\langle i|h|j\rangle
  = \int d\mathbf{x}_1\, \chi_i^*(\mathbf{x}_1)\hat{h}(\mathbf{r}_1)\chi_j(\mathbf{x}_1)
\end{equation}
\begin{equation}
 \langle ij|kl\rangle=\langle \chi_i\chi_j|\chi_k\chi_l\rangle
=\int d\mathbf{x}_1 d\mathbf{x}_2\,\chi_i^*(\mathbf{x}_1)\chi_j^*(\mathbf{x}_2)
r_{12}^{-1} \chi_k(\mathbf{x}_1)\chi_l(\mathbf{x}_2)
\end{equation}

\begin{equation}
  \left[ij||kl\right]=\langle ij||kl\rangle
  = \int d\mathbf{x}_1 d\mathbf{x}_2\,
     \chi_i^*(\mathbf{x}_1)\chi_j^*(\mathbf{x}_2)\,
     r_{12}^{-1}\,
     \bigl[ \chi_k(\mathbf{x}_1)\chi_l(\mathbf{x}_2)
          - \chi_l(\mathbf{x}_1)\chi_k(\mathbf{x}_2) \bigr]
\end{equation}
空间轨道
\begin{equation}
 \left( i|h|j\right)=h_{ij}=\int d\mathbf{r}_1\,\psi_i^*(\mathbf{r}_1)
 h(\mathbf{r}_1)\psi_j(\mathbf{r}_1)
\end{equation}
\begin{equation}
  \left( ij|kl\right)=\int d\mathbf{r}_1 d\mathbf{r}_2\,\psi_i^*(\mathbf{r}_1)
  \psi_j(\mathbf{r}_1)r_{12}^{-1}\psi_k^*(\mathbf{r}_2)\psi_l(\mathbf{r}_2)
\end{equation}
\begin{equation}
  \text{库仑积分 } J_{ij}=\left(ii|jj\right)
\end{equation}
\begin{equation}
  \text{交换积分 } K_{ij}=\left(ij|ji\right)
\end{equation}
\subsubsection{General Rules for Matrix Elements}
本节介绍一套求矩阵元的规则，下一节展示推导.量子化学有两类算符.第一类是单电子算符之和
\begin{equation}
  \hat{O}_1=\sum_{i=1}^{N}h\left(i\right)
\end{equation}
第二类是两电子算符之和
\begin{equation}
  \hat{O}_2=\sum_{i=1}^{N}\sum_{j>i}v\left(i,j\right)=\sum_{i<j}v\left(i,j\right)
\end{equation}
(1)单电子算符在行列式间的矩阵元有三种情况
case1:$|K\rangle=|\dots mn \dots\rangle$
\begin{equation}
  \langle K|\hat{O}_1|K\rangle=\sum_{m}^{N}\langle m|h|m \rangle
\end{equation}
case2:$|K\rangle=|\dots mn \dots\rangle$ 
$|L\rangle=|\dots pn \dots\rangle$ 
\begin{equation}
  \langle K|\hat{O}_1|L\rangle=\sum_{m}^{N}\langle m|h|p \rangle
\end{equation}
case3:$|K\rangle=|\dots mn \dots\rangle$
$|L\rangle=|\dots pq \dots\rangle$
\begin{equation}
  \langle K|\hat{O}_1|L\rangle=0
\end{equation}
(2)双电子算符在行列式间的矩阵元有三种情况
case1:$|K\rangle=|\dots mn \dots\rangle$
\begin{equation}
  \langle K|\hat{O}_2|K\rangle=\frac{1}{2}\sum_{m}^{N}\sum_{n}^{N}\langle mn||mn \rangle
\end{equation}
case2:$|K\rangle=|\dots mn \dots\rangle$ 
$|L\rangle=|\dots pn \dots\rangle$ 
\begin{equation}
  \langle K|\hat{O}_2|L\rangle=\sum_{n}^{N}\langle mn||pn \rangle
\end{equation}
case3:$|K\rangle=|\dots mn \dots\rangle$
$|L\rangle=|\dots pq \dots\rangle$
\begin{equation}
  \langle K|\hat{O}_2|L\rangle=\langle mn||pq \rangle
\end{equation}
根据上列式子可以快速写出单个行列式$|K\rangle$的能量
\begin{equation}
  \langle K|\hat{H}|K\rangle=
  \sum_{m}^{N}\langle m|h|m\rangle +\frac{1}{2}\sum_{m}^{N}\sum_{n}^{N}\langle mn||mn\rangle
  \label{eq:determinant-energy}
\end{equation}
其中
\begin{equation}
  h(i)=-\frac{1}{2}\nabla_i^2-\sum_{A}\frac{Z_A}{r_{iA}}
\end{equation}
对于$H_2$极小基下，$|\Phi_0\rangle=|\chi_1\chi_2\rangle$
\begin{equation}
  \begin{aligned}
  E_0&=\langle 1|h|1\rangle+\langle2|h|2\rangle+\langle12||12\rangle \\
     &=\langle 1|h|1\rangle+\langle2|h|2\rangle+\langle12|12\rangle-\langle12|21\rangle
  \end{aligned}
\end{equation}
\subsubsection{Derivation of the Rules for Matrix Elements}
略
\subsubsection{Transition from Spin Orbitals to Spatial Orbitals}
尽管在前面章节中使用自旋轨道可以简化量子化学理论设计的代数运算和各种记号，但是在运算中，
自旋函数$\alpha$和$\beta$必须被积分掉.才能将自旋轨道约化为可数值运算的空间轨道及其积分

我们以$H_2$极小基的Hartree-Fock基态为例
\begin{equation}
  E_0=\langle \chi_1|h|\chi_1\rangle + \langle \chi_2|h|\chi_2\rangle + \langle 12||12\rangle
\end{equation}  
可以得到
\begin{equation}
  E_0=\langle \psi_1|h|\psi_1\rangle+\langle \bar{\psi_1}|h|\bar{\psi_1}\rangle+\langle \psi_1\bar{\psi_1}||\psi_1\bar{\psi_1}\rangle
\end{equation}
对于单电子积分
\begin{equation}
 \langle \bar{\psi_1}|h|\bar{\psi_1}\rangle=\int d\mathbf{r}d\omega_1\psi_1^*(\mathbf{r})\beta^*(\omega_1)h(\mathbf{r})\psi_1(\mathbf{r})\beta(\omega_1)
\end{equation}
注意到$\langle\beta|\beta\rangle$,那么就有
\begin{equation}
  \langle \bar{\psi_1}|h|\bar{\psi_1}\rangle=\int d\mathbf{r}\psi_1^*(\mathbf{r})h(\mathbf{r})\psi_1(\mathbf{r})=\langle \psi_1|h|\psi_1\rangle
\end{equation}
所以$E_0$中单电子部分的能量贡献是$2\langle \psi_1|h|\psi_1\rangle$
对于$\langle \psi_1\bar{\psi_1}||\psi_1\bar{\psi_1}\rangle$有
\begin{equation}
  \begin{aligned}
    \langle \psi_1\bar{\psi_1}||\psi_1\bar{\psi_1}\rangle
    &= \langle \psi_1\bar{\psi_1}|\psi_1\bar{\psi_1}\rangle - \langle \psi_1\bar{\psi_1}|\bar{\psi_1}\psi_1\rangle \\
    &= \iint d\mathbf{r}_1 d\mathbf{r}_2\, \psi_1^*(\mathbf{r}_1)\beta^*(\omega_1)\psi_1^*(\mathbf{r}_2)\alpha^*(\omega_2) r_{12}^{-1} \\
    &\qquad\times \bigl[ \psi_1(\mathbf{r}_1)\beta(\omega_1)\psi_1(\mathbf{r}_2)\alpha(\omega_2) - \psi_1(\mathbf{r}_1)\alpha(\omega_1)\psi_1(\mathbf{r}_2)\beta(\omega_2) \bigr] 
  \end{aligned}
\end{equation}
对于括号中的第一项，由于$\beta$和$\alpha$的正交性，该矩阵元会因子化为空间积分与自旋内积的乘积
\begin{equation}
  \beta(\omega_1)|\beta(\omega_1)\rangle =1 
\end{equation}
\begin{equation}
  \alpha(\omega_2)|\alpha(\omega_2)\rangle =1
\end{equation}
所以
\begin{equation}
  \iint r_{12}^{-1}\psi_1(1)\psi_1(1)\psi_1(2)\psi_1(2)d\tau_1d\tau_2=\langle \psi_1\psi_1|\psi_1\psi_1\rangle
\end{equation}
对于第二项，自旋内积是$\beta(\omega_1)|\alpha(\omega_1)\rangle=0$，所以整个项为0
因此$\langle \psi_1\bar{\psi_1}||\psi_1\bar{\psi_1}\rangle=\langle \psi_1\psi_1|\psi_1\psi_1\rangle$
也就是
\begin{equation}
  \begin{aligned}
  E_0&=2\langle\psi_1|h|\psi_1\rangle+\langle\psi_1\psi_1|\psi_1\psi_1\rangle \\
     &=2\langle 1|h|1\rangle+\langle 11|11\rangle
  \end{aligned}
\end{equation}

把上面的结论推广到有偶数电子的N电子体系中，可以写出N电子系统的闭壳层限制性Hartree-Fock波函数
\begin{equation}
  \Psi_0=|\psi_1\bar{\psi_1}\psi_2\bar{\psi_2}\cdots\psi_a\bar{\psi_a}\cdots\psi_N\bar{\psi_N}|
\end{equation}
积分可得
\begin{equation}
  E_0=2\sum_i\langle i|h|i\rangle+\sum_{i,j}(2J_{ij}-K_{ij})
  \label{eq:rhf-energy}
\end{equation}
\subsubsection{Coulomb and Exchange Integrals}
式\eqref{eq:rhf-energy}中单电子项
\begin{equation}
  \langle a|h|a\rangle=h_{aa}=\int d\mathbf{r}_1\,\psi_a^*(\mathbf{r}_1)\left(-\frac{1}{2}\nabla_1^2-\sum_{A}\frac{Z_A}{r_{1A}}\right)\psi_a(\mathbf{r}_1)
\end{equation}
$h_{aa}$就是波函数$\psi_a(\mathbf{r}_1)$所对应电子的平均动能加上核对它的吸引能量

式\eqref{eq:rhf-energy}中双电子项
\begin{equation}
  J_{ij}=\langle ij|ij\rangle =
  \int d\mathbf{r}_1 d\mathbf{r}_2\,|\psi_i(\mathbf{r}_1)|^2 r_{12}^{-1} |\psi_j(\mathbf{r}_2)|^2
\end{equation}
就是函数$|\psi_a(\mathbf{r}_1)|^2$电子云与$|\psi_b(\mathbf{r}_2)|^2$电子云的库仑相互作用的静电能，并称之为库仑积分.
而对于积分
\begin{equation}
  K_{ij}=\langle ij|ji\rangle =  \int d\mathbf{r}_1 d\mathbf{r}_2\,\psi_i^*(\mathbf{r}_1)\psi_j^*(\mathbf{r}_2) r_{12}^{-1} \psi_j(\mathbf{r}_1)\psi_i(\mathbf{r}_2)
\end{equation}
没有经典解释对应，称为交换积分.库仑积分和交换积分的值都是正的.下面我们证明行列式对应的能量中的交换积分是
由交换相关 (exchange correlation) 作用引起的，也即在单行列式近似下, 同自旋电子的运动相互关联.

\vspace{1em}
\noindent\rule{\linewidth}{0.4pt}

极小基$H_2$的full CI矩阵为
\begin{equation}
\hat{H}=\begin{bmatrix}
  2h_{11}+J_{11}&K_{12} \\
  K_{12}&2h_{22}+J_{22}
\end{bmatrix}
\label{eq:h2-full-ci}
\end{equation}
\textbf{证明：} full CI 基函数取 $|\Psi_0\rangle=|1\bar{1}\rangle$ 和
$|\Psi_{1\bar{1}}^{2\bar{2}}\rangle=|2\bar{2}\rangle$

\medskip
\noindent\textbf{(1) 对角元 $\langle 1\bar{1}|\hat{H}|1\bar{1}\rangle$}

单电子部分：
\[
\langle 1\bar{1}|\hat{O}_1|1\bar{1}\rangle
= \langle 1|h|1\rangle + \langle \bar{1}|h|\bar{1}\rangle
= h_{11}+h_{11}=2h_{11}
\]
双电子部分
\[
\langle 1\bar{1}|\hat{O}_2|1\bar{1}\rangle
= \frac12(J_{11}+J_{11})=J_{11}
\]
因此
\[
\langle 1\bar{1}|\hat{H}|1\bar{1}\rangle = 2h_{11}+J_{11}
\]
同理
\[
\langle 2\bar{2}|\hat{H}|2\bar{2}\rangle = 2h_{22}+J_{22}
\]

\medskip
\noindent\textbf{(2) 非对角元 $\langle 1\bar{1}|\hat{H}|2\bar{2}\rangle$}

两行列式相差两个自旋轨道（$1\to2$, $\bar{1}\to\bar{2}$），
单电子部分为零：
\[
\langle 1\bar{1}|\hat{O}_1|2\bar{2}\rangle = 0
\]

双电子部分：
\[
\langle 1\bar{1}|\hat{O}_2|2\bar{2}\rangle
= \langle 1\bar{1}||2\bar{2}\rangle
= \langle 1\bar{1}|2\bar{2}\rangle - \langle 1\bar{1}|\bar{2}2\rangle
\]

积分掉自旋：
\[
\langle 1\bar{1}|2\bar{2}\rangle
= (12|12)\,\langle\alpha|\alpha\rangle\langle\beta|\beta\rangle
= (12|12)
\]
\[
\langle 1\bar{1}|\bar{2}2\rangle
= (12|12)\,\langle\alpha|\beta\rangle\langle\beta|\alpha\rangle = 0
\]

由于空间轨道为实函数，$(12|12)=(12|21)=K_{12}$。故
\[
\langle 1\bar{1}|\hat{H}|2\bar{2}\rangle = K_{12}
\]

同理 $\langle 2\bar{2}|\hat{H}|1\bar{1}\rangle = K_{12}$
（哈密顿矩阵为实对称矩阵）。

\smallskip
综上即得 \eqref{eq:h2-full-ci}。


\noindent\rule{\linewidth}{0.4pt}
\vspace{1em}
之前已经说到, 在一个两电子自旋平行的系统$|\bar{\psi_1}\bar{\psi_2}\rangle$中, 
在同一点同时找到这两个电子的概率为零.但若这两个电子自旋相反,
即系统此时的波函数为$|\psi_1\bar{\psi_2}\rangle$, 则此概率不为零. 那么一个合理的推测
就是电子间的库仑排斥会使态$|\bar{\psi_1}\bar{\psi_2}\rangle$对应的能量比
$|\psi_1\bar{\psi_2}\rangle$的能量低. 利用式\eqref{eq:determinant-energy}，
$|\psi_1\bar{\psi_2}\rangle$(即为$|\uparrow\downarrow\rangle$)的能量就是
\begin{equation}
  E_{\uparrow\downarrow}=\langle\psi_1\bar{\psi_2}|\hat{H}|\psi_1\bar{\psi_2}\rangle
=h_{11}+h_{22}+J_{12}
\end{equation}
而$|\bar{\psi_1}\bar{\psi_2}\rangle$(即为$|\downarrow\downarrow\rangle$或$|\uparrow\uparrow\rangle$)的能量为
\begin{equation}
  E_{\downarrow\downarrow}=E_{\uparrow\uparrow}
=h_{11}+h_{22}+J_{12}-K_{12}
\end{equation}
由于$K_{12}$是正值，前者能量小于后者.因此, Slater 行列式所对应的能量的表达式中交换积分的出现,
 实际上反 映了这样一个事实:即使在单行列式近似下, 有相同自旋的两个电子的运动也是相互关联着的(相关,correlated).
 在\ref{sec:HartreeProduct}里提到, 
 Hartree 乘积中没有交换积分出现. 也就不能满足费米子的要求. 
 所以电子之间并没有相互规避, 电子运动是无关的(没有相关), 因此不能正确反映体系的性质.
\subsubsection{Pseudo-Classical Interpretation of Determinantal Energies}
先讨论能量中的单电子贡献.有如下结论: 空间轨道$\psi_i$上的一个电子 (无论自旋如何) 会在总能量中贡献一项$h_{ii}$.
双电子贡献里一对电子既可以平行自旋也可以相反.若自旋相反
\begin{equation}
  \langle ij||ij\rangle=J_{ij}
\end{equation}
若自旋平行
\begin{equation}
  \langle ij||ij\rangle=J_{ij}-K_{ij}
\end{equation}
因此我们有如下论断:空间轨道$\psi_i$、$\psi_j$中的每一对电子 (无论其自旋) 都对能量贡献一项$J_{ij}$.
空间轨道$\psi_i$、$\psi_j$中的一对自旋平行电子对能量贡献一项$-K_{ij}$. 行列式对应的总能就是这些项之和.

按这种方式来表述能量时需记住, 自旋平行的电子之间的交换作用只不过是一种方便的表述方式,
是为了将由单行 列式所描述的系统的能量以简单的方式表达出来, 它本身并无真实物理意义.
若真的论及两电子间真实的物理相互作用, 我们只能说这种相互作用由哈密顿量中的库仑排斥项 ($r_{ij}^{-1}$) 所描述,
而与自旋无关.
\subsection{Second Quantization}
反对称原理是量子力学中的一个公理,此前我们通过Slater行列式及其线性组合来保证波函数满足这个原理.而是否能找到其他办法满足反对称原理呢？
二次量子化就是这样一种处理多电子体系的有效框架.借助二次量子化, 我们可以把重心从 N-电子波函数本身上移开, 
转而关注之前研究过的单双电子积分$\langle i|h|j\rangle$,$\langle i||j\rangle$ .本书其余章节并未
广泛地使用二次量子化, 因此本节可以选读.

略
\subsection{Spin-Adapted Configuration}
本节会 更详细地讨论自旋, 并讨论多电子系统的自旋态.
\subsubsection{Spin Operators}
单个粒子的自旋角动量算符是矢量算符
\begin{equation}
  \tilde{s}=\tilde{s}_x i+\tilde{s}_y j+\tilde{s}_z k
\end{equation}
可以用$s^2$及其某个分量的共同本征函数集合作为单粒子自旋态的完备集, 一般将分量选为$s_z$
\begin{equation}
  s^2=|s,m_s\rangle=s(s+1)|s,m_s\rangle
  ,\quad
  s_z|s,m_s\rangle=m_s|s,m_s\rangle
\end{equation}
上式中s是描述总自旋的一个量子数，s可以取$0,\frac{1}{2},1,\frac{3}{2},$并取相应值$m_s=-s,-s+1,\dots,s$对于电子而言,s总为1/2,因此不需在$s$做标记
\begin{equation}
  s_z|\alpha\rangle=\frac{1}{2}|\alpha\rangle
  ,\quad s_z|\beta\rangle=-\frac{1}{2}|\beta\rangle
\end{equation}
有常用恒等式
$$\begin{aligned}
  s_x|\alpha\rangle&=\frac{1}{2}|\beta\rangle,
  s_x|\beta\rangle=\frac{1}{2}|\alpha\rangle ,\\
  s_y|\alpha\rangle&=\frac{i}{2}|\beta\rangle,
  s_y|\beta\rangle=-\frac{i}{2}|\alpha\rangle,\\
  s_z|\alpha\rangle&=\frac{1}{2}|\alpha\rangle,
  s_z|\beta\rangle=-\frac{1}{2}|\beta\rangle.
\end{aligned}$$

此外还有升/降算符, 常用来处理自旋
\begin{equation}
  s_+=s_x+is_y,
  s_-=s_x-is_y\end{equation}
这两个算符可以升降$s_z$的值一个单位:
\begin{equation}
  s_+|\beta\rangle=|\alpha\rangle,
  s_+|\alpha\rangle=0,
  s_-|\alpha\rangle=|\beta\rangle,
  s_-|\beta\rangle=0
\end{equation}
也就是说,$s_+$作用于$\beta$自旋时产生$\alpha$自旋.

多电子体系的总自旋算符就是每个电子自旋算符的矢量和
\begin{equation}
  S=\sum_{i=1}^N s(i)
\end{equation}
$S^2$与$S_Z$哈密顿量对易，那么哈密顿量的精确本征函数也就是这两个自旋算字的本征函数
\begin{equation}
  S^2|\Psi\rangle=S(S+1)|\Psi\rangle
  ,\quad S_Z|\Psi\rangle=M_S|\Psi\rangle
\end{equation}
式中$S$、$M_S$是N电子多重态$|\Phi\rangle$总自旋量子数及其z方向上的分量.总自旋磁量子数,取值为$-S,-S+1,\dots,S$
一般而言,多电子波函数的自旋多重度(angular multiplicity) 是所对应的这个值$2S+1$.我们分别把这些态称作单重态、双重态
、三重态等.
\subsubsection{Restricted Determinants and Spin-Adapted Configurations}
最简单的闭壳层行列式就是极小基$H_2$ 的 Hartree-Fock 基态波函数
\begin{equation}
  |\Psi_0\rangle=|\psi_1\bar{\psi}_1|=[\psi_1(1)\psi_1(2)]2^{-\frac{1}{2}}(\alpha_1\beta_2-\alpha_2\beta_1)
\end{equation}
上式最后展开了行列式.这个波函数的自旋部分就是一双电子系统的单重自旋 函数. 同样, 双激发态$|\psi_{1\bar{1}}^{2\bar{2}}\rangle$也是单重态.

现来看开壳层限制性行列式.对于极小基$H_2$四个单激发行列式中的两个
\begin{equation}
  | \Psi_{1}^{\bar{2}} \rangle
  =|\bar{21}\rangle=-2^{-\frac{1}{2}}[\psi_1(1)\psi_2(2)-\psi_2(1)\psi_1(2)]\beta_1\beta_2
\end{equation}
\begin{equation}
  | \Psi_{\bar{1}}^{2} \rangle
  =|12\rangle=2^{-\frac{1}{2}}[\psi_1(1)\psi_2(2)-\psi_2(1)\psi_1(2)]\alpha_1\alpha_2
\end{equation}
是$S^2$的本征函数且本征值为2，均为三重态.但剩下的两个
\begin{equation}
  | \Psi_{\bar{1}}^{\bar{2}} \rangle
  =|\bar{21}\rangle=-2^{-\frac{1}{2}}[\psi_1(1)\psi_2(2)+\psi_2(1)\psi_1(2)]\beta_1\alpha_2
\end{equation}
\begin{equation}
  | \Psi_1^{2} \rangle
  =|12\rangle=2^{-\frac{1}{2}}[\psi_1(1)\psi_2(2)+\psi_2(1)\psi_1(2)]\beta_1\alpha_2
\end{equation}$S^2$的本征函数,而是一个单重态和三重态的线性组合.很明显,在很多情况下,人们需要构建是$S^2$本征函数的组态.对于双电子系统,态很简单,$S^2$归一化的本征函数为
\begin{equation}
  {}^{3}\Psi_1^{2}=2^{-\frac{1}{2}}(|\bar{21}\rangle-|12\rangle)
\end{equation}

\section{The Hartree-Fock Approximation}
本章中来详细介绍 Hartree-Fock 理论以及从头的 Hartree-Fock 计算.
 此章篇幅较长, 这也反映了 Hartree-Fock 在量子化学中的重要地位. 
 Hartree-Fock 不仅就其本身十分重要, 它也是其他考虑了电子相关效应的更精确近似的起点.
\subsection{The Hartree-Fock Equations}
本节是对方程推导结果的总结，具体推导细节放在下一节讨论.

我们需要找到一组自旋轨道，使其构成的单行列式
\begin{equation}
  |\Psi_0 \rangle=|\chi_1\chi_2\cdots\chi_N\rangle
\end{equation}
成为该形式下N电子系统基态的最佳近似.根据变分原理，这种轨道就是使得下式能量最小的轨道
\begin{equation}
  E_0=\langle \Psi_0|\hat{H}|\Psi_0\rangle
  =\sum_a^{N}\langle a|h|a\rangle + \frac{1}{2}\sum_{ab}^{N}\langle ab||ab\rangle
\end{equation}
在下一小节中可以推导唯一的限制就是要求他们相互正交
\begin{equation}
  \langle a|b\rangle=\delta_{ab}
\end{equation}
可以得到一个刻画最优自旋轨道的方程，Hartree-Fock积分微分方程
\begin{equation}
  \begin{split}
    h(1)\chi_a(1)&+\sum_{b\neq a}\biggl[\int dx_2\,|\chi_b(2)|^2 r_{12}^{-1}\biggr]\chi_a(1) -\sum_{b\neq a}\biggl[\int dx_2\,\chi_b^*(2)\chi_a(2) r_{12}^{-1}\biggr]\chi_b(1) \\
    &=\epsilon_a\chi_a(1)
  \end{split}
  \label{eq:hf-diff}
\end{equation}
其中
\begin{equation}
  h(1)=-\frac{1}{2}\nabla_1^2-\sum_{A}\frac{Z_A}{r_{1A}}
\end{equation}
\subsubsection{Coulomb and Exchange Operators}
设想电子2占据$\chi_b$，Hartree-Fock理论中，电子1感受到电子2在空间瞬时位置产生的双电子势$_{12}^{-1}$被替换为单电子势，同时考虑占据$d\mathbf{x}_2$
的概率权重$d\mathbf{x}_2|\chi_b(2)|^2$.由此可以定义库仑算符
\begin{equation}
  \mathcal{J}_b(1)=\int dx_2|\chi_b(2)|^{2}r_{12}^{-1}
  \label{eq:coulombo}
\end{equation}
同时引入交换算符，根据其对自旋轨道的作用定义
\begin{equation}
  \mathcal{K}_b(1)\chi_a(1)=\biggl[\int dx_2\chi_b^*(2)r_{12}^{-1}\chi_a(2)\biggr]\chi_b(1)
  \label{eq:exchangeo}
\end{equation}
交换项来自单行列式的反对称性质, 它的形式比较特殊, 而且没有如库 伦项那样简单的经典解释. 但仍可将 Hartree-Fock 方程\eqref{eq:hf-diff}写为
\begin{equation}
  \begin{split}
    &\biggl[ h(1)+\sum_{b\neq a}\mathcal{J}_b(1)-\sum_{b\neq a}\mathcal{K}_b(1)\biggr]\chi_a(1)=\epsilon_a\chi_a(1)
  \end{split}
\end{equation}
这里交换算符是非定域算符，因为在空间一点$\mathbf{x}_1$上并不存在一个简单的势能$\mathcal{J}_b(\mathbf{x}_1)$.实际上
这个算符的作用结果依赖于$\chi_a$在空间上分布的情况.因此, 交换算符的作用结果是一个函数, 而不是一个数值.也无法画出交换势的等高线图，但库仑势就可以.

对于在$\chi_a$上的电子而言, 库仑算符和交换算符$\mathcal{J}_b$,$\mathcal{K}_b$ 的期望值 就是前一章提过的库仑和交换积分
\begin{equation}
  \langle a|\mathcal{J}_b|a\rangle=J_{ab}
\end{equation}
\begin{equation}  
  \langle a|\mathcal{K}_b|a\rangle=K_{ab}
\end{equation}
\subsubsection{The Fock Operators}
对于式\ref{eq:coulombo}与\ref{eq:exchangeo},有如下结论
\begin{equation}
  [\mathcal{J}_a(1)-\mathcal{K}_a(1)]\chi_a(1)=0
\end{equation}
将这一项加入方程\eqref{eq:hf-diff}中，有:
\begin{equation}
  f(1)\chi_a(1)=\epsilon_a\chi_a(1)
\end{equation}
也即
\begin{equation}
  f|\chi_a\rangle=\epsilon_a|\chi_a\rangle
  \label{hf}
\end{equation}
其中Fock算符定义为
\begin{equation}
  f(1)=h(1)+\sum_b^{N}[\mathcal{J}_b(1)-\mathcal{K}_b(1)]
\end{equation}
上式后两项称为有效单电子势算符，亦称Hartree-Fock势
\begin{equation}
  \nu^{HF}(1)=\sum_b( \mathcal{J}_b(1)-\mathcal{K}_b(1))
\end{equation}
与单个电子算符类似，可以定义Fock算符在给定自旋轨道$\chi_a$上的期望值.且本征值$\epsilon_a$恰为轨道能量.
并且,HF算符显然是一组完整正交的轨道.而电子总能量的表达式:
\begin{equation}
  E_0=\sum_a\langle a|h|a\rangle+\frac{1}{2}\sum_{ab}\langle ab||ab\rangle
\end{equation}
Fock算符的总和=Fock算符的本征值之和-双电子积分之和:
\begin{equation}
  \sum_a\langle a|f|a\rangle=\sum_a\epsilon_a=\sum_a\langle a|h|a\rangle+\sum_{ab}\langle ab||ab\rangle
\end{equation}
显然，Fock算符的总和即总Hamiltonian在单行列式下的期望值.

为方便，Fock算符可用$\mathcal{P}_{12}$重新写为
\begin{equation}
  f(1)=h(1)+\sum_b\int dx_2\chi_b^*(2)r_{12}^{-1}(1-\mathcal{P}_{12})\chi_b(2)
\end{equation}
\subsection{Derivation of the Hartree-Fock Equations}
本节来推导 Hartree-Fock 方程的一般自旋轨道形式, 也即, 将单 Slater 行列式的 能量极小化以得到式\ref{hf}.
\subsubsection{Functional Variation}
变分的目的就是寻找一个合适的$\Phi$使得$E[\Phi]$达到最小
\begin{equation}
  \delta E=0
\end{equation}
这个条件保证了E对于$\Phi$的任何变分都处于驻态(stationary).通常情况下，驻点就是最小值.
给定一个线性变分尝试波函数
\begin{equation}
  |\Phi\rangle=\sum_{i=1}^{N}c_i|\Psi_i\rangle
\end{equation}
希望得到能量极小值
\begin{equation}
  E=\langle \Phi|\hat{H}|\Phi\rangle=\sum_{ij}c_i^*c_j\langle\Psi_i|\hat{H}|\Psi_j\rangle
\end{equation}
并且保持归一化
\begin{equation}
\langle\Phi|\Phi\rangle-1=\sum_{ij}c_i^*c_j\langle\Psi_i|\Psi_j\rangle-1=0
\end{equation}
利用Lagrange不定乘子法，我们可以将如下的泛函对$c_i$极小化
\begin{equation}
  \mathcal{L}=\sum_{ij}c_i^*c_j\langle\Psi_i|\hat{H}|\Psi_j\rangle-E\biggl(\sum_{ij}c_i^*c_j\langle\Psi_i|\Psi_j\rangle-1\biggr)
\end{equation}
下面令$\mathcal{L}$的一节变分为0
\begin{equation}
  \begin{aligned}
  \delta\mathcal{L}&=\sum_{ij}\delta c_i^*c_j\langle\Psi_i|\hat{H}|\Psi_j\rangle-E\sum_{ij}\delta c_i^*c_j\langle\Psi_i|\Psi_j\rangle \\
  &\quad +\sum_{ij}c_i^*\delta c_j\langle\Psi_i|\hat{H}|\Psi_j\rangle-E\sum_{ij}c_i^*\delta c_j\langle\Psi_i|\Psi_j\rangle
  \end{aligned}
\end{equation}
两两合并，后两项交换指标利用厄米性可得
\begin{equation}
  \sum_{i}\delta c_i^*\big[\sum_j( H_{ij}-ES_{ij}c_{j})\big]+c.c. =0
\end{equation}
那么对任意的$\delta c_i^*$,都必有
\begin{equation}
\sum_j( H_{ij}-ES_{ij})c_j =0
\end{equation}
或写为
\begin{equation}
  \mathbf{Hc}=E\mathbf{Sc}
\end{equation}
我们将用它来导出Hartree-Fock方程
\subsubsection{minimization of the Energy of a Single Determinant}
给定一个单行列式 $|\Psi_0\rangle=|\chi_1\chi_2\dots\chi_N\rangle$, 能量$E_0\langle \Psi_0|\hat{H}|\Psi_0\rangle$是自旋轨道${\chi_a}$的泛函.
 为导出 Hartree-Fock方程需要关于自旋轨道进行极小化, 而且要满足自旋轨道正交归一的约束
\begin{equation}
  \langle \chi_a|\chi_b\rangle-\delta_{ab}=0
\end{equation}
考虑$\mathcal{L}[\{\chi_a\}]$作为自旋轨道的泛函
\begin{equation}
  \mathcal{L}[\{\chi_a\}]=E_0[\{\chi_a\}]-\sum_{ab}\epsilon_{ab}(\langle \chi_a|\chi_b\rangle-\delta_{ab})
\end{equation}
其中$E_0$是单行列式$|\Psi_0\rangle$的期望值
\begin{equation}
  E_0[\{\chi_a\}]=\sum_a^{N}\langle a|h|a\rangle + \frac{1}{2}\sum_{ab}^{N}\langle ab||ab\rangle
\end{equation}
仿照上一节,$\mathcal{L}$的一阶变分为0
\begin{equation}
  \begin{aligned}
  \delta\mathcal{L} &= \sum_{a=1}^{N}\int d\mathbf{x}_1\,\delta\chi_a^*(1)
    \Big[ h(1)\chi_a(1) + \sum_{b=1}^{N}\bigl(\mathcal{J}_b(1)-\mathcal{K}_b(1)\bigr)\chi_a(1)-\sum_{b=1}^N\epsilon_{ba}\chi_b(1) \Big]+c.c. \\
  &= \sum_{a=1}^{N}\int d\mathbf{x}_1\,\delta\chi_a^*(1)
    \Big[ f(1)\chi_a(1) - \sum_{b=1}^N\epsilon_{ba}\chi_b(1) \Big]
  \end{aligned}
\end{equation}
又有$\delta\chi_a^*(1)$是任意的，上式方括号内对所有a都等于0
\begin{equation}
  f(1)\chi_a(1) = \sum_{b=1}^N\epsilon_{ba}\chi_b(1)
\end{equation}
也即
\begin{equation}
  f|\chi_a\rangle = \sum_{b=1}^N\epsilon_{ba}|\chi_b\rangle
  \label{eq:Fock-like}
\end{equation}
这个方程并非标准的本征值形式，所以这个式子看起来有些奇怪.但我们可以通过自旋轨道之间的酉变换导出正则形式
\subsubsection{The Canonical Hartree-Fock Equations}
对于一组新自旋轨道${\chi_a'}$由旧的自旋轨道${\chi_a}$通过酉变换而来
\begin{equation}
  |\chi_a'\rangle=\sum_b U_{ab}|\chi_a\rangle
\end{equation}
酉变换实际上就是满足以下关系的变换
\begin{equation}
  \mathbf{U^\dagger}=\mathbf{U^{-1}}
\end{equation}
它能保持 (被变换基的) 正交归一性.现考虑一经过酉变换的波函数，有
\begin{equation}
  |\Psi_0'\rangle=det(\mathbf{U})|\Psi_0\rangle
\end{equation}
因为
\begin{align*}
  \mathbf{U^\dagger}\mathbf{U} &= 1 \\
  \det(\mathbf{U^\dagger}\mathbf{U}) &= |\det(\mathbf{U})|^2 = 1
\end{align*}
所以
\begin{equation}
  \det(\mathbf{U})=e^{i\theta}
\end{equation}

所以变换后行列式与原行列式仅差一个相因子.若$\mathbf{U}$是酉矩阵，则仅能取±1.
在单行列式 波函数下, 所有期望值在自旋轨道的任意酉变换下都不变. 因此使总能量取驻值的
自旋轨道集合并不是唯一的, 而且某一组自旋轨道并不比其他轨道更有物理意义. 举个例 子,
 定域自旋轨道并不比离域自旋轨道更“物理”.我们利用单行列式在自旋轨道酉变换下不变的性质,
 将式\ref{eq:Fock-like}改写为 关于一组特定轨道的本征值方程. 但首先需确定以上
 酉变换对 Fock 算符和 Lagrange 乘子矩阵的影响.

 Fock 算符中依赖于自旋轨道的部分是库仑和交换算符. 经过变换后的库仑算符之和为:
\begin{align*}
  \sum_a\mathcal{J}_a'(1) &= \sum_a \int d\mathbf{x}_2\, \chi_a^{\prime*}(2) r_{12}^{-1} \chi_a'(2) \\
  &= \sum_{bc}\Bigl[\sum_a U_{ba}^* U_{ca}\Bigr] \int d\mathbf{x}_2\, \chi_b^*(2) r_{12}^{-1} \chi_c(2)
\end{align*}
注意到
\begin{equation}
  \sum_a U_{ba}^*U_{ca}=\delta_{bc}
\end{equation}
那么可以得到
\begin{equation}
  \sum_a\mathcal{J}_a'(1)=\sum_b\mathcal{J}_b(1)
\end{equation}
因此库仑算符在自旋轨道的酉变换下不变. 以相同的方式易证交换算符之和在自旋轨道的任意酉变换下不变, 因而 Fock 算符也不变
\begin{equation}
  f'(1)=f(1)
\end{equation}
现在来确定酉变换对 Lagrange 乘子的作用. 将式\ref{eq:Fock-like}左乘$\langle\chi_c|$就可证明 Lagrange 乘子是 Fock 算符的矩阵元
\begin{equation}
  \epsilon_{ca}=\langle\chi_c|f|\chi_a\rangle
\end{equation}
\begin{equation}
  \begin{aligned}
    \epsilon_{ab}' &= \int d\mathbf{x}_1\,\chi_a^{\prime*}(1)f(1)\chi_b'(1) \\
    &= \sum_{cd} U_{ca}^*U_{db} \int d\mathbf{x}_1\,\chi_c^*(1)f(1)\chi_d(1) \\
    &= (\mathbf{U^\dagger\epsilon U})_{ab}
  \end{aligned}
\end{equation}
对于$\epsilon$是一个厄米矩阵.因此总能找到一个合适的酉矩阵$\mathbf{U}$使得变换$\epsilon'=\mathbf{U^\dagger\epsilon U}$把
$\epsilon$对角化.
并且它存在且唯一. 那么一定 有一组自旋轨道 ${\chi_a′}$ 使得 Lagrange 乘子矩阵为对角化,Hartree-Fock 方程写为
\begin{equation}
  f|\chi_a\rangle=\epsilon_a|\chi_a\rangle
\end{equation}
这组解叫作正则自旋轨道.它们一般是离域化的并且构成该分子点群下的不可约表 示的基, 也即, 
它们会有对应于该分子对称性的特定对称性质, 或者等价地, 对应于 Fock 算符的对称性. 
一旦得到这组正则自旋轨道, 我们就能对其进行酉变换得到无数组等价 的轨道. 
实际上, 有许多准则可用来寻找酉变换以使轨道经变换后变为定域的、更符合 对化学键的直觉的轨道.
\subsection{Interpretation of Solutions to the Hartree-Fock Equations}
为求解 Hartree-Fock 方程, 需要引入基组并求解对应的矩阵方程. 但该本征值方 程及其解也有不依赖于基组的一些性质, 此时正好来讨论它们.
\subsubsection{Orbital Energies and Koopmans' Theorem}
对于 N-电子体系, 将行列式 $|\chi_1\chi_2\dots\chi_N\rangle$的能量极小化后就得到一组关 于 N 个被占自旋轨道 ${\chi_a}$
 的本征值方程$f|\chi_a\rangle=\epsilon_a|\chi_a\rangle$. Fock 算符对这些被占轨道 有泛函的依赖, 
 但若这些被占轨道已知, 那么 Fock 算符就变为一个正常的厄米算符, 有 无穷个本征函数
\begin{equation}
  f|\chi_j\rangle=\epsilon_j|\chi_j\rangle\quad j=1,2,\dots,\infty
  \label{eq:fock}
\end{equation}
式\ref{eq:fock}的每个解$|\chi_i\rangle$都有一个自旋轨道能量$\epsilon_j$.对于 N-电子体系, 在 Hartree-Fock 理论中 N 个最低自旋轨道能量对应 N 个被占自旋轨道,
其他的解称为虚自旋轨道. 它们可以由 Fock 算符对角化得到, 这些解构成了一个无穷集合.
用$\langle\chi_i|$左乘式\ref{eq:fock}可以证明Fock 算符在这组自旋轨道下的的矩阵表示是对角 的, 其对角元就是轨道能量
\begin{equation}
  \langle\chi_i|f|\chi_j\rangle=\epsilon_j\delta_{ij}
\end{equation}
轨道能量可以写作
\begin{equation}
  \epsilon_j=\langle j|h|j\rangle+\sum_{b=1}^N\bigl[\langle jb||jb\rangle\bigr]
\end{equation}
具体地, 根据以上式子我们有
\begin{equation}
  \epsilon_a=\langle a|h|a\rangle+\sum_{b=1}^N\bigl(\langle ab||ab\rangle\bigr)
\end{equation}
\begin{equation}
  \epsilon_r=\langle r|h|r\rangle+\sum_{b=1}^N\bigl(\langle rb||rb\rangle\bigr)
\end{equation}
若我们将基态$|\Psi_0\rangle$中的 N 个电子的能量$\epsilon_a$简单相加, 会得到
\begin{equation}
  \sum_a^N\epsilon_a=\sum_a^N\langle a|h|a\rangle+\sum_{ab}^N\langle ab||ab\rangle
\end{equation}
与单行列式的总能量$E_0$的表达式
\begin{equation}
  E_0=\sum_a^N\langle a|h|a\rangle+\frac{1}{2}\sum_{ab}^N\langle ab||ab\rangle
\end{equation}
相比, 前者多了$\frac{1}{2}\sum_{ab}^N\langle ab||ab\rangle$.因此, 轨道能量的总和并不等于单行列式的总能量. 
但它们之间的差距仅为双电子积分项的一半,这是因为每两个自旋轨道上的两电子间相互作用被计入了两次.

若总能并非轨道能量之和，现考虑其他物理意义.电子的湮灭和产生可以提供答案.
Koopmans 定理 给定一个 $N$-电子 Hartree-Fock 单行列式 $|^{N}\Psi_0\rangle$,
其中占据和虚自旋轨道能量分别为 $\epsilon_a$, $\epsilon_r$, 则从 $\chi_a$ 中挪去一个电子,
保持其余轨道不变, 会产生一个 $(N-1)$-电子单行列式 $|^{N-1}\Psi_a\rangle$, 
这个过程对应的电离能是 $-\epsilon_a$. 在虚轨道 $\chi_r$ 上增加一个电子, 保持自旋轨道不变, 
生成一个 $(N+1)$ 电子单行列式 $|^{N+1}\Psi^r\rangle$, 此过程的电子亲和能为 $-\epsilon_r$.

Koopmans 定理给出一种计算近似的电离能和电子亲和能的办法. 这种“冻结轨道”近似是假设 $(N \pm 1)$-电子态
中的自旋轨道与 $N$-电子态中的一样. 这个近似忽略了 $(N \pm 1)$-电子态中自旋轨道的弛豫, 
就是说, $|^{N}\Psi_0\rangle$ 中的自旋轨道并非 $|^{N-1}\Psi_a\rangle$ / $|^{N+1}\Psi^r\rangle$ 的最优自旋轨道. 
如果分别再用另外的 Hartree-Fock 计算优化这两个态, 则 $^{N-1}E_a$, $^{N+1}E_r$ 的能量要减少.
因此 Koopmans 定理忽略了弛豫会使算出的电离能和电子亲和能过正. 当然除此之外,
对波函数作单行列式近似也会产生误差, 即相关效应, 通过 Hartree-Fock 之外的办法可以计算, 
可对 Koopmans 定理的结果作进一步修正. 实际上, 当系统内电子数目越大时, 相关能也会越大. 
因此相关效应会抹掉电离能的一部分弛豫误差, 而在电子亲和能的弛豫误差上再增加误差. 一般来说, 
Koopmans 电离能是实验电离能值合理的一阶近似, 本章后面会讨论大量的此类计算.
可惜 Koopmans 电子亲和能会较差. 很多中性分子增加一个电子可形成稳定的负离子. 
但中性分子的 Hartree-Fock 计算常给出正的虚轨道能量. 电子亲和能是非常难以计算的量, 
要比电离能难算得多
\subsubsection{Brillouin's Theorem}
正如上一章讨论的, 由集 合 ${\chi_i}$ 可以构造相当多的行列式. 知道了 Fock 算符的形式,
现在正好可以证明一个关 于这些行列式的子集的定理. 该子集就是所有单激发行列式$|\Psi_r^a\rangle$,
此即——将 $|\Psi_0\rangle$
中 $\chi_a$ 换为 $\chi_r$ 即得. 在精确基态的多行列式表示中, 我们可能会预想, 这些单行列式将给出领头阶修正
\begin{equation}
  |\Psi_0\rangle=|\Psi_0\rangle+\sum_{ra}c_a^r|\Psi_a^r\rangle
\end{equation}
若在修正中仅考虑单激发行列式, 那么 $c_a^r$ 可以由变分原理得到——将 $\{\Psi_0\}, \{\Psi_a^r\}$ 基下的哈密顿矩阵对角化即可. 现在我们来研究包含单激发态的矩阵本征值问题:
\begin{equation}
  \begin{bmatrix}
    \langle\Psi_0|\hat{H}|\Psi_0\rangle & \langle\Psi_0|\hat{H}|\Psi_a^r\rangle \\
    \langle\Psi_a^r|\hat{H}|\Psi_0\rangle & \langle\Psi_a^r|\hat{H}|\Psi_a^r\rangle
  \end{bmatrix}  \begin{bmatrix}
    c_0 \\ c_a^r
  \end{bmatrix}
  =\epsilon
  \begin{bmatrix}
    c_0 \\ c_a^r
  \end{bmatrix}
  \label{eq:smatrix}
\end{equation}
两个态的混合体现在非对角元$\langle\Psi_0|\hat{H}|\Psi_a^r\rangle$上,这个矩阵元可以用我们介绍过的行列 式间矩阵元计算规则求出
\begin{equation}
  \langle\Psi_0|\hat{H}|\Psi_a^r\rangle=\langle a|h|r\rangle+\sum_b^N\langle ab||rb\rangle
\end{equation}
对于方程右边，回想Fock算符的矩阵元可以得到
\begin{equation}
  \langle a|f|r\rangle=\langle\Psi_0|\hat{H}|\Psi_a^r\rangle
\end{equation}
由于正则 Hartree-Fock 方程 $f|\chi_i\rangle=\epsilon_i|\chi_i\rangle$ 给出正交本征函数，对于
被占轨道 $a$ 与虚轨道 $r$（$a\neq r$）有 $\langle a|f|r\rangle = \epsilon_r\langle a|r\rangle = 0$.
因此 $\langle\Psi_0|\hat{H}|\Psi_a^r\rangle = 0$，即单激发行列式与 Hartree-Fock 基态间的矩阵元为零。
所以求解 Hartree-Fock 本征值方程等价于确保 $|\Psi_0\rangle$ 不和任何单激发行列式发生混合.
式\ref{eq:smatrix}中能量最低的解为
\begin{equation}
  \begin{pmatrix}
  E_0 & 0 \\
  0 & \langle\Psi_a^r|\hat{H}|\Psi_a^r\rangle
  \end{pmatrix}
  \begin{pmatrix}
    1 \\
    0
  \end{pmatrix}
  =E_0
  \begin{pmatrix}
    1 \\
    0
  \end{pmatrix}
\end{equation}
在这种意义下, Hartree-Fock 基态可以说是“稳定”的, 因为即使将单激发行列式混入, 结果也不能进一步改进. 
Brillouin 定理单激发行列式 $|\Psi_a^r\rangle$ 不会直接与参考态即 Hartree-Fock 行列式发生作用, 
也即 $\langle\Psi_0|\hat{H}|\Psi_a^r\rangle = 0$.
\subsubsection{The Hartree-Fock Hamiltonian}
目前为止, 我们对 Hartree-Fock 近似的理解就是:Hamiltonian 是精确的, 
而波函 数用单 Slater 行列式来近似. 本节中,为了第六章微扰论中的需要, 
提前以另外一种等 价但不同的视角研究 Hartree-Fock 理论, 此时我们主要关注的是 Hamiltonian 本身.

是否存在某个近似的 N-电子 Hamiltonian,其本征值方程和我们之前精确求解过的完全相同?或者说,
 是否存在某近似的 Hamiltonian, 使$|\Psi_0\rangle$就是其精确本征函数?答案 是存在.
  Hartree-Fock Hamiltonian 定义如下
  \begin{equation}
    \hat{H}_0 =\sum_i^N\hat{f}(i)
  \end{equation}
上式对应本征值不是Hartree-Fock能量$E_0$，而是轨道能量之和$\sum_{a}\epsilon_a$.
实际我们可以证明由 Fock 算符的本征函数 $\{\chi_i\}$ 组成的任意行列式都是 $H_0$ 的本征函数,
 本征值等于行列式中自旋轨道能量之和.

\end{document}